\section{Method}\label{sec:method}

\paragraph{Setting.}
Each missense variant has a predictor score $s$, oriented so that higher means more damaging, and an assay label
$Y\in\{0,1\}$, with $Y=1$ for loss of function (LoF). A \emph{residue} is a protein position of one gene; the
substitutions at a residue share structural context and form one cluster. A gene's assayed missense variants are split
by residue into a calibration part $\mathcal A$ and a certification part $\mathcal B$; the variants to be called (for
example VUS) form the target. We write $q(s)\in[0,1]$ for an increasing map of the score that puts the three
predictors on one scale (here Platt scaling \cite{Platt1999} fitted to confidently labelled ClinVar variants, which the
fixed-cut baselines of \cref{sec:data} also use). Because $q$ is strictly increasing, every \textsc{CertiVar} call is the
same on $q$ as on $s$, and no ClinVar label affects any call.

\paragraph{Per-class thresholds.}
The nonconformity of label 1 is $r_1=1-q$ and of label 0 is $r_0=q$. From each residue in $\mathcal A$ we keep one
variant chosen uniformly at random. For class $k$, let $n_k$ be the number of kept variants with label $k$ and
$r^{(k)}_{(1)}\le\dots\le r^{(k)}_{(n_k)}$ their sorted nonconformities. The threshold is $\hat t_k=r^{(k)}_{(j_k)}$
($\infty$ if $j_k>n_k$), where $j_k$ is the in-expectation rank or the probably approximately correct (PAC) rank at
confidence level $\delta_{\mathrm P}$ ($j_k=n_k+1$ if no $j$ qualifies):
\begin{equation}\label{eq:rank}
j_k=\bigl\lceil (n_k+1)(1-\alpha_k)\bigr\rceil \;\text{(in expectation)}\quad\text{or}\quad
j_k=\min\bigl\{\,j:\ F_{\mathrm{Beta}(j,\,n_k-j+1)}(1-\alpha_k)\le\delta_{\mathrm P}\,\bigr\}\;\text{(PAC)}.
\end{equation}
The prediction set is $C(q)=\{k: r_k\le \hat t_k\}$; \textsc{CertiVar} calls LoF if $C=\{1\}$, functional if
$C=\{0\}$, and abstains otherwise. In effect it places two cut-offs on the score, $\hat t_0$ and $1-\hat t_1$, each an
order statistic of one class: a variant below both is called functional, one above both is called LoF, and one between
them is not called (\cref{fig:overview}). A functional call on a LoF variant requires $1\notin C$, so per-class coverage
bounds both error rates.

\paragraph{Certified likelihood ratios.}
The ACMG/AMP guidelines \cite{Richards2015} grade evidence in qualitative categories; in their Bayesian adaptation,
supporting, moderate, strong and very strong evidence towards pathogenicity correspond to a likelihood ratio above 2.1,
4.3, 18.7 and 350, and supporting, moderate and strong evidence towards benignity to one below 0.48, 0.23 and 0.053
\cite{Tavtigian2018,Brnich2020}. We reuse these values as a reporting scale, test levels up to strong, and write, for example,
\emph{LoF-supporting} for LoF calls whose likelihood ratio exceeds 2.1. Our likelihood ratios are defined with respect to
the assay outcome $Y$ and averaged over all variants that receive a call, as OddsPath is for a functional assay
\cite{Brnich2020}; they are not the local, score-specific ratios that ClinGen uses for the computational criteria
PP3/BP4 \cite{Pejaver2022}. A level thus describes the call class as a whole, need not hold for each call, and must not
be combined with the same assay's functional evidence (PS3/BS3) or with PP3/BP4 from the same predictor. For thresholds
$t$, the call rates are $\lambda_k(t)=\Pr(\text{call LoF}\mid Y=k)$ and $\phi_k(t)=\Pr(\text{call functional}\mid Y=k)$,
and $\mathrm{LR}^+=\lambda_1/\lambda_0$, $\mathrm{LR}^-=\phi_1/\phi_0$.

For each $\alpha$ on a grid $G$ we compute thresholds on $\mathcal A$ ($\alpha_0=\alpha_1=\alpha$), call one variant per
residue of $\mathcal B$, and bound the four rates by one-sided Clopper--Pearson bounds \cite{ClopperPearson1934} at
level $\epsilon=\delta/(2|G|)$: for $x$ calls among $n$ variants of a class, the lower bound is
$F^{-1}_{\mathrm{Beta}(x,n-x+1)}(\epsilon)$ ($0$ if $x=0$) and the upper bound $F^{-1}_{\mathrm{Beta}(x+1,n-x)}(1-\epsilon)$
($1$ if $x=n$, including $n=0$). Writing $\underline{\lambda}_1,\overline{\lambda}_0,\overline{\phi}_1,\underline{\phi}_0$
for these bounds, the certified bounds are $L^+(\alpha)=\underline{\lambda}_1/\overline{\lambda}_0$ and
$U^-(\alpha)=\overline{\phi}_1/\underline{\phi}_0$ ($U^-=\infty$, never meeting a level, if $\underline{\phi}_0=0$;
$\overline{\lambda}_0>0$ always; for point estimates, $x/0=\infty$ for $x>0$, and $0/0$ is undefined and never
selected). We write $\hat t(\alpha)$ for the pair $(\hat t_0,\hat t_1)$ at $\alpha$ and call the share of the kept variants of
$\mathcal B$ that receive a call of a given type its \emph{yield}. For each evidence level, \cref{alg:certivar} returns
the $\alpha$ with the largest yield among those whose certified bound meets the level. Because $\hat t_0$ and $\hat t_1$ both decrease weakly in $\alpha$, no variant can be
called LoF at one $\alpha$ and functional at another (\cref{sec:theory}). Line~7 requires, for functional calls, $\alpha$
at least as large as the LoF choice of the same strength, so that the two call types of one strength come from an
ordered pair of thresholds; this is a reporting convention, not needed for validity, and can only lower functional
yields. The \emph{fixed-sequence} variant replaces the Bonferroni level $\delta/(2|G|)$ by $\delta/(2L)$, where $L$ is
the number of evidence levels tested for a call type, tests the grid in an order fixed in advance, and stops at the first
$\alpha$ that does not meet the level (Theorem~\ref{thm:fs}). We test LoF levels in increasing and functional levels in
decreasing $\alpha$, because few variants are called functional at small $\alpha$. All thresholds and bounds cost
$O(|G|\,n\log n)$ for $n$ residues, about 5\,ms per gene on a laptop CPU.

\paragraph{Checking class-conditional invariance.}
The guarantees transfer to a target set whose LoF share differs from the calibration data only if, within each class,
target scores follow the calibration law (Assumption~\ref{as:shift} below). This implies that the target score law is a
mixture $\pi F_1+(1-\pi)F_0$ of the two class laws for some $\pi\in[0,1]$. We test this necessary condition on one
target variant per residue with
$D=\min_{\pi}\sup_x\bigl|\hat F_T(x)-\pi\hat F_1(x)-(1-\pi)\hat F_0(x)\bigr|$, where $\hat F_T$, $\hat F_1$ and
$\hat F_0$ are the empirical distribution functions of the target scores and of the two classes in $\mathcal A$, and
$\pi$ runs over a grid of step 0.005. The null distribution of $D$ is approximated by a bootstrap that redraws both
calibration classes and a target sample of the same size from the fitted mixture (300 resamples). A rejection shows
that the target is not a re-weighting of the calibration classes, so the guarantees need not transfer, and we do not
report calls for such a target set; acceptance does not establish Assumption~\ref{as:shift}.

\begin{algorithm}[t]
\caption{\textsc{CertiVar} for one gene and one evidence level (Bonferroni version)}\label{alg:certivar}
\begin{algorithmic}[1]
\Require scores $q$; assay labels split by residue into $\mathcal A,\mathcal B$; grid $G$; $\delta$; levels $\tau$
(LoF) and $\tau'$ (functional)
\State keep one variant per residue of $\mathcal A$ and of $\mathcal B$
\For{$\alpha\in G$}
\State thresholds $\hat t_0,\hat t_1$ on the kept variants of $\mathcal A$ by \eqref{eq:rank}; call the kept variants of
$\mathcal B$
\State $L^+(\alpha)\gets \underline{\lambda}_1/\overline{\lambda}_0$, $U^-(\alpha)\gets \overline{\phi}_1/\underline{\phi}_0$
(Clopper--Pearson, level $\delta/(2|G|)$ each)
\EndFor
\State $\hat\alpha^+\gets\arg\max_{\alpha\in G}\{\text{share called LoF on }\mathcal B : L^+(\alpha)>\tau\}$
\State $\hat\alpha^-\gets\arg\max_{\alpha\in G}\{\text{share called functional on }\mathcal B : U^-(\alpha)<\tau',\ \alpha\ge\hat\alpha^+\}$
(over all $\alpha\in G$ if $\hat\alpha^+$ does not exist)
\State \textbf{output} LoF calls at $\hat\alpha^+$ and functional calls at $\hat\alpha^-$ with their certified levels (no
calls of a type whose set is empty); abstain otherwise
\end{algorithmic}
\end{algorithm}

\subsection{Guarantees}\label{sec:theory}

Let $Q$ be the law of a variant obtained by drawing a residue from the gene's population of assayed residues and then
one of its variants uniformly at random (binary-labelled, non-splice variants outside the ClinVar labelled set). Every
rate and likelihood ratio below is taken under $Q$, so it weights residues equally.

\begin{assumption}[Clustered sampling]\label{as:cluster}
Residues in $\mathcal A$, in $\mathcal B$ and of a test variant are independent draws from one residue population;
within a residue, the kept or tested variant is chosen uniformly at random; $q$ does not depend on these data.
\end{assumption}

\begin{assumption}[Class-conditional invariance]\label{as:shift}
Target variants have, for $k\in\{0,1\}$, the same law of $q$ given $Y=k$ as a test variant under
Assumption~\ref{as:cluster}; their LoF share $\pi=\Pr(Y=1)$ may differ.
\end{assumption}

\begin{proposition}[Per-class coverage]\label{prop:mondrian}
Under Assumption~\ref{as:cluster}, for a fixed $\alpha_k$, with the in-expectation rank and conditionally on $n_k$,
$1-\alpha_k\le\Pr(k\in C\mid Y=k,n_k)\le 1-\alpha_k+1/(n_k+1)$, the upper bound holding when there are no ties.
\end{proposition}

\begin{proposition}[High-probability per-class coverage]\label{prop:pac}
Under Assumption~\ref{as:cluster}, for a fixed $\alpha_k$ and with the PAC rank, $\Pr(k\in C\mid Y=k,\hat t_k)\ge1-\alpha_k$
with probability at least $1-\delta_{\mathrm P}$ over the calibration data for each $k$, and for both classes with
probability at least $1-2\delta_{\mathrm P}$.
\end{proposition}

\begin{proposition}[Coverage after selection]\label{prop:select}
Under Assumption~\ref{as:cluster}, for any $\hat\alpha\in G$ chosen from data independent of the test variant
(for example by \cref{alg:certivar}), the in-expectation rank gives $\Pr(k\in C\mid Y=k)\ge1-\max G$; the PAC rank at
level $\delta_{\mathrm P}/|G|$ gives $\Pr(k\in C\mid Y=k,\hat t_k)\ge1-\hat\alpha$ for each $k$ with probability at least
$1-\delta_{\mathrm P}$.
\end{proposition}

\begin{theorem}[Certified likelihood ratios]\label{thm:lr}
Under Assumption~\ref{as:cluster}, conditionally on $\mathcal A$ and with probability at least $1-\delta$ over
$\mathcal B$, $\mathrm{LR}^+(\hat t(\alpha))\ge L^+(\alpha)$ for every $\alpha\in G$ simultaneously; the same holds for
$\mathrm{LR}^-\le U^-$ with a second $\delta$. Because the event does not depend on the level, each call type returned
by \cref{alg:certivar} reaches its certified level at every evidence level reported, with probability at least
$1-\delta$ (both call types with probability at least $1-2\delta$), whichever $\alpha$ is selected.
\end{theorem}

\begin{theorem}[Fixed-sequence certification]\label{thm:fs}
Under Assumption~\ref{as:cluster}, fix a level $\tau$, a confidence $\delta'$ and an order of $G$ chosen in advance.
Reject $H_\alpha:\mathrm{LR}^+(\hat t(\alpha))\le\tau$ when $\underline{\lambda}_1/\overline{\lambda}_0>\tau$ with both
endpoints at level $\delta'/2$, stop at the first $\alpha$ not rejected, and select the rejected $\alpha$ with the
largest yield; the rejected set is a prefix of the order. (a) Conditionally on $\mathcal A$, the selected LoF calls have
$\mathrm{LR}^+>\tau$ with probability at least $1-\delta'$. (b) If each of $L$ levels is tested in this way with
$\delta'=\delta/L$, every level claimed for LoF calls holds simultaneously with probability at least $1-\delta$. The same
holds for functional calls with $\overline{\phi}_1/\underline{\phi}_0<\tau'$, and for both call types together with
probability at least $1-2\delta$.
\end{theorem}

\begin{corollary}[No prior needed]\label{cor:shift}
Under Assumptions~\ref{as:cluster} and~\ref{as:shift}, for target variants drawn residue first as under $Q$,
Propositions~\ref{prop:mondrian}--\ref{prop:select} and Theorems~\ref{thm:lr}--\ref{thm:fs} hold whatever their LoF
share $\pi\in(0,1)$: for a fixed
$\alpha$, LoF variants are called functional at rate at most $\alpha_1$, and the certified likelihood ratios with
respect to the assay outcome need no prior. For any $\alpha^+,\alpha^-\in G$, the LoF calls at $\alpha^+$ and the
functional calls at $\alpha^-$ are disjoint.
\end{corollary}

Propositions~\ref{prop:mondrian}--\ref{prop:pac} are the class-conditional and training-conditional conformal
guarantees \cite{Vovk2005,Sadinle2019,Vovk2013,Tong2018} under $Q$; Theorems~\ref{thm:lr}--\ref{thm:fs} apply the union
bound of risk control \cite{Bates2021rcps,Geifman2017} and the fixed-sequence argument of Learn-then-Test
\cite{Angelopoulos2025ltt} to a \emph{ratio} of class-conditional rates of an abstaining rule. Monotonicity of both ranks
in $\alpha$ makes coverage at a selected $\alpha$ at least that at $\max G$ and keeps the call types disjoint. Proofs are
in \cref{app:proofs}.

\paragraph{Scope.}
After selection, per-class coverage is guaranteed only at $1-\max G$; Theorem~\ref{thm:fs} certifies levels, not
bounds. Assumption~\ref{as:cluster} treats residues as independent draws (our random splits draw without replacement;
the oracle design does not), and a contiguous region need not be exchangeable with the rest; Assumption~\ref{as:shift}
holds only approximately for VUS (\cref{sec:shift}) and not for variants whose assay result would be intermediate.

